\begin{figure}[h!]
\centering
\begin{tikzpicture}[
  font=\small,
  >=Stealth,
  stage/.style={
    draw=black!65,
    rounded corners=2pt,
    fill=white,
    align=center,
    minimum height=10mm,
    text width=19mm,
    inner sep=3pt
  },
  profile/.style={
    stage,
    fill=violet!7,
    text width=34mm,
    minimum height=12mm
  },
  public/.style={
    stage,
    fill=blue!7,
    text width=24mm
  },
  flow/.style={->, line width=.75pt},
  config/.style={->, dashed, gray!75, line width=.6pt}
]

% 1st colour-con diag
    % initial 2 dots
    \fill (0,0) circle[radius=2pt];
    \fill (0.5,0) circle[radius=2pt];

    % brackets
        % 1 
        \draw (1,-0.5) -- (1,0.5);
        \draw (1,-0.5) -- (1.5,-0.5);
        \draw (1.5,-0.5) -- (1.5,0.5);
        % 2
        \draw (2,-0.5) -- (2,0.5);
        \draw (2,-0.5) -- (2.5,-0.5);
        \draw (2.5,-0.5) -- (2.5,0.5);
        % 3
        \draw (3,-0.5) -- (3,0.5);
        \draw (3,-0.5) -- (3.5,-0.5);
        \draw (3.5,-0.5) -- (3.5,0.5);
        % 4
        \draw (4,-0.5) -- (4,0.5);
        \draw (4,-0.5) -- (4.5,-0.5);
        \draw (4.5,-0.5) -- (4.5,0.5);

    % final 2 dots
    \fill (5,0) circle[radius=2pt];
    \fill (5.5,0) circle[radius=2pt];

    % text
    \node[above] at (1.75,0.5) {$i$};
    \node[above] at (2.75,0.5) {$j$};
    \node[above] at (3.75,0.5) {$k$};

% arrow
\draw[->] (2.75,-1) -- (2.75,-1.5);

% 2nd colour-con diag
    % initial 2 dots
    \fill (0,-2.5) circle[radius=2pt];
    \fill (0.5,-2.5) circle[radius=2pt];
    % brackets
        % 1 
        \draw (1,-3) -- (1,-2);
        \draw (1,-3) -- (1.5,-3);
        \draw (1.5,-3) -- (1.5,-2);
        % 2
        \draw (2,-3) -- (2,-2);
        \draw (2,-3) -- (3.5,-3);
        \draw (3.5,-3) -- (3.5,-2);
        % 3
        \draw (4,-3) -- (4,-2);
        \draw (4,-3) -- (4.5,-3);
        \draw (4.5,-3) -- (4.5,-2);

    % final 2 dots
    \fill (5,-2.5) circle[radius=2pt];
    \fill (5.5,-2.5) circle[radius=2pt];

    % text
    \node[above] at (1.75,-2) {$I$};
    \node[above] at (3.75,-2) {$K$};

% Amplitude-factorisation diagram, reproducing Fig. 3 of Gehrmann-De
% Ridder, Gehrmann and Glover (hep-ph/0505111) -- the (m+1)-parton
% amplitude blob (with generic spectator legs 1,...,m+1 either side of
% i,j,k) factorising into the reduced m-parton amplitude blob times, in
% square brackets, the tree antenna function X^0_{ijk} itself drawn as
% a cut amplitude-squared diagram. The paper draws this left-to-right;
% here it is rotated to a vertical top-to-bottom flow (top diagram,
% arrow down, bottom diagram) so it sits compactly next to the
% colour-connection panel on the left, which uses the same top/arrow/
% bottom layout.
\node[right] at (6.9,-1.25) {
\begin{tikzpicture}[scale=0.72,every node/.style={scale=0.72}]

    % ----- top: full (m+1)-parton amplitude blob -----
    % shifted right so this block's own centre lines up with the centre
    % of the (much wider) bottom row -- blob + bracket -- below it
    \begin{scope}[shift={(2.3,0)}]
        \draw (-1.0,0) -- (0,0);
        \fill (0,0) circle[radius=2.0mm];
        \draw (0,0) -- (55:1.4); \node[above] at (55:1.5) {$1$};
        \fill (45:0.85) circle[radius=1.6pt];
        \fill (34:0.85) circle[radius=1.6pt];
        \draw (0,0) -- (25:1.4); \node[right] at (25:1.5) {$i$};
        \draw (0,0) -- (0:1.4);  \node[right] at (0:1.5)  {$j$};
        \draw (0,0) -- (-25:1.4); \node[right] at (-25:1.5) {$k$};
        \fill (-34:0.85) circle[radius=1.6pt];
        \fill (-45:0.85) circle[radius=1.6pt];
        \draw (0,0) -- (-55:1.4); \node[below] at (-55:1.55) {$m+1$};
    \end{scope}

    % ----- vertical arrow, centred under the top block above -----
    % kept at its original short length -- the extra room between the
    % top and bottom blocks is left as plain empty space around it,
    % not absorbed into a longer arrow
    \draw[->] (2.3,-2.0) -- (2.3,-2.7);

    % ----- bottom: reduced m-parton amplitude blob, and, aligned with -----
    % ----- it, the bracket giving the tree antenna function X^0_{ijk} -----
    % ----- as a cut amplitude-squared diagram                        -----
    \begin{scope}[shift={(0,-5.6)}]
        \draw (-1.0,0) -- (0,0);
        \fill (0,0) circle[radius=2.0mm];
        \draw (0,0) -- (55:1.4); \node[above] at (55:1.5) {$1$};
        \fill (45:0.85) circle[radius=1.6pt];
        \fill (34:0.85) circle[radius=1.6pt];
        \draw (0,0) -- (20:1.4); \node[right] at (20:1.5) {$I$};
        \draw (0,0) -- (-20:1.4); \node[right] at (-20:1.5) {$K$};
        \fill (-34:0.85) circle[radius=1.6pt];
        \fill (-45:0.85) circle[radius=1.6pt];
        \draw (0,0) -- (-55:1.4); \node[below] at (-55:1.55) {$m+1$};

        % the bracket is deliberately bigger than the reduced-blob diagram
        % to its left -- more headroom around the two antenna sub-diagrams
        % it contains, so the "k"/"I" legs don't crowd the divider line.
        % Extra clearance added: between the reduced blob and the bracket's
        % left edge, between the bracket edges and the mini-diagrams they
        % contain, and between the divider line's ends and the bracket
        % sides.
        \draw (2.95,2.0) -- (2.7,2.0) -- (2.7,-2.0) -- (2.95,-2.0);
        \draw (5.97,2.0) -- (6.22,2.0) -- (6.22,-2.0) -- (5.97,-2.0);
        \draw (3.25,0) -- (5.67,0);

        \begin{scope}[shift={(4.15,1.0)}]
            \draw[dash pattern=on 1.2pt off 1pt] (-0.9,0) -- (0,0);
            \fill (0,0) circle[radius=2.0mm];
            \draw (0,0) -- (25:1.35); \node[right] at (25:1.45) {$i$};
            \draw (0,0) -- (0:1.35);  \node[right] at (0:1.45)  {$j$};
            \draw (0,0) -- (-25:1.35); \node[right] at (-25:1.45) {$k$};
        \end{scope}

        \begin{scope}[shift={(4.15,-1.0)}]
            \draw[dash pattern=on 1.2pt off 1pt] (-0.9,0) -- (0,0);
            \fill (0,0) circle[radius=2.0mm];
            \draw (0,0) -- (25:1.35); \node[right] at (25:1.45) {$I$};
            \draw (0,0) -- (-25:1.35); \node[right] at (-25:1.45) {$K$};
        \end{scope}
    \end{scope}

\end{tikzpicture}
    };

\end{tikzpicture}
\caption{\centering The colour connections in between parent and daughter partons before and after an emission, accompanied also by a schematic Feynman diagram using the momenta labels.}
\label{fig:colourConnectionNLO}
\end{figure}
